\section{Example of a RenderScript kernel writing and usage}
\label{app:ex2}
\noindent In this example we compute the luminance histogram of an image, but we want to do it in two separate steps:
\begin{itemize}
\item in the first step the image is splitted in stripes and for each one the histogram is computed.
\item the second step sums the histograms obtained from the previously step.
\end{itemize}
For implementing the program in the way described above we need to write two different kernel functions.
\\\\
This is the code for the RenderScript runtime framework, which resides in our .rs file:

\begin{lstlisting}[frame=single]

// The histogram script requires a number of buffers

// The source and destination images.
rs_allocation gSrc;
rs_allocation gDest;

// A buffer for the intermediate sums
// Integer values, [256][steps] in size
rs_allocation gSums;

// Final sum buffer, Integer, [256] cells in size
rs_allocation gSum;

// The height and width of the image
int gWidth;
int gHeight;

// The step, which is the number of lines processed by each thread
int gStep;

// The number of steps in total, roughly height / step
int gSteps;

// The start of the first kernel
void __attribute__((kernel)) pass1(int in, uint x, uint y) {
  // Note that x and y will indicate the coordinates of the pixel being processed
  // This kernel will be run on a range of x = [0] and y = [0 to steps]
  // Clear our output buffer for this thread.
  for (int i=0; i < 256; i++) {
    // Set the value at i,y to 0
    rsSetElementAt_int(gSums, 0, i, y);

  }

  // Iterate over our image
  for (int i= 0; i < gStep; i++) {
    int py= y*gStep + i;	// Compute the y coordinate in the image
    if (py>= gHeight) return;	// Might be out of range if this is the last step

    // Walk one scanline
    for (int px=0; px < gWidth; px++) {
      // Get the pixel and convert to luminance
      uchar4 c = rsGetElementAt_uchar4(gSrc, px, py);
      int lum = (77 * c.r + 150 * c.g + 29 * c.b) >> 8;

      // Add one to the bucked for this luminance value
      int old = rsGetElementAt_int(gSums, lum, y);
      rsSetElementAt_int(gSums, old+1, lum, y);
    }
  }
}

// This kernel is run on the Sum allocation, so its called once
// for each of the 256 levels
// Note, this is a 1D kernel
int __attribute__((kernel)) pass2(uint x) {
  int sum = 0;

  // For this level, add in the sum from each of the
  // separate partial sums
  for (int i=0; i < gSteps; i++) {
    sum += rsGetElementAt_int(gSums, x, i);
  }

  // Return the sum for this level. Since this is a kernel
  // return value, it will automatically be placed in the allocation.
  return sum;
}

// This is an invokable function. It will be called single threaded
void rescale() {
  // Find our largest bucket value
  int maxv = 0;
  for (int i=0; i < 256; i++) {
    maxv = max(maxv, rsGetElementAt_int(gSum, i));
  }

  // Compute the rescale value to to convert bucket values into bar heights.
  float overMax = (1.f / maxv) * gHeight;

  for (int i=0; i < 256; i++) {
    int t = rsGetElementAt_int(gSum, i);
    t = gHeight - (overMax * rsGetElementAt_int(gSum, i));
    t = max(0, t);
    rsSetElementAt_int(gSum, t, i);
  }
}


\end{lstlisting}

\vspace{4ex}
This is the code we have to write in the Java application that allows to launch our scripts:

\begin{lstlisting}[frame=single]

// Create and load the script
mScript = new ScriptC_histogram(mRS);

// Setup the globals
mScript.set_gWidth(width);
...

// Create the [256][steps] buffer of integers for the partial sums
Type.Builder tb = new Type.Builder(mRS, Element.I32(mRS));
tb.setX(256).setY(steps);
Type t = tb.create();
mSums = Allocation.createTyped(mRS, t);

// Create the 1D [256] buffer for the final Sums
mSum = Allocation.createSized(mRS, Element.I32(mRS), 256);

// Set the buffers for the script
mScript.set_gSums(mSums);

// This first pass should be clipped because we want [step] threads
// not [256][step] threads

// To do this we have the ability to clip our launch
// using a LaunchOptions class
Script.LaunchOptions lo = new Script.LaunchOptions();

// Set the range in the X dimension to be 0 to 1
// Note: the end is exclusive so this says only run X value 0
lo.setX(0, 1);

// Run the kernel with our launch options
// This will spawn one thread per Y coordinate, each with X=0
mScript.forEach_pass1(mSums, lo);

// Once the first pass is complete, we need to add up our partial sums

// The pass2 launch is unclipped.
// It spawns one thread per cell in the mSum buffer
// for a total of 256 threads.
mScript.forEach_pass2(mSum);

// Finally, we call our rescale function
mScript.invoke_rescale();


\end{lstlisting}
\vspace{4ex}
This is an example of an image with its luminance histogram:

\begin{figurehere}
 \centering
 \includegraphics[width=0.5\textwidth]{./pictures/bitmap}
 \label{fig:bitmap}
\end{figurehere}
\begin{figurehere}
 \centering
 \includegraphics[width=0.5\textwidth]{./pictures/lumhist}
 \label{fig:lumhist}
\end{figurehere}
